% Options for packages loaded elsewhere
\PassOptionsToPackage{unicode}{hyperref}
\PassOptionsToPackage{hyphens}{url}
\documentclass[
  a4paper,
]{article}
\usepackage{xcolor}
\usepackage[margin=25mm]{geometry}
\usepackage{amsmath,amssymb}
\setcounter{secnumdepth}{-\maxdimen} % remove section numbering
\usepackage{iftex}
\ifPDFTeX
  \usepackage[T1]{fontenc}
  \usepackage[utf8]{inputenc}
  \usepackage{textcomp} % provide euro and other symbols
\else % if luatex or xetex
  \usepackage{unicode-math} % this also loads fontspec
  \defaultfontfeatures{Scale=MatchLowercase}
  \defaultfontfeatures[\rmfamily]{Ligatures=TeX,Scale=1}
\fi
\usepackage{lmodern}
\ifPDFTeX\else
  % xetex/luatex font selection
\fi
% Use upquote if available, for straight quotes in verbatim environments
\IfFileExists{upquote.sty}{\usepackage{upquote}}{}
\IfFileExists{microtype.sty}{% use microtype if available
  \usepackage[]{microtype}
  \UseMicrotypeSet[protrusion]{basicmath} % disable protrusion for tt fonts
}{}
\makeatletter
\@ifundefined{KOMAClassName}{% if non-KOMA class
  \IfFileExists{parskip.sty}{%
    \usepackage{parskip}
  }{% else
    \setlength{\parindent}{0pt}
    \setlength{\parskip}{6pt plus 2pt minus 1pt}}
}{% if KOMA class
  \KOMAoptions{parskip=half}}
\makeatother
\setlength{\emergencystretch}{3em} % prevent overfull lines
\providecommand{\tightlist}{%
  \setlength{\itemsep}{0pt}\setlength{\parskip}{0pt}}
\usepackage{newunicodechar}
\newunicodechar{↔}{\ensuremath{\leftrightarrow}}
\newunicodechar{→}{\ensuremath{\rightarrow}}
\newunicodechar{←}{\ensuremath{\leftarrow}}
\newunicodechar{≥}{\ensuremath{\ge}}
\newunicodechar{≤}{\ensuremath{\le}}
\newunicodechar{×}{\ensuremath{\times}}
\newunicodechar{≪}{\ensuremath{\ll}}
\newunicodechar{≫}{\ensuremath{\gg}}
\newunicodechar{∈}{\ensuremath{\in}}
\newunicodechar{∞}{\ensuremath{\infty}}
\newunicodechar{±}{\ensuremath{\pm}}
\newunicodechar{−}{\ensuremath{-}}
\newunicodechar{≈}{\ensuremath{\approx}}
\newunicodechar{≠}{\ensuremath{\ne}}
\newunicodechar{∝}{\ensuremath{\propto}}
\newunicodechar{√}{\ensuremath{\sqrt{\phantom{x}}}}
\newunicodechar{⊕}{\ensuremath{\oplus}}
\newunicodechar{α}{\ensuremath{\alpha}}
\newunicodechar{β}{\ensuremath{\beta}}
\newunicodechar{θ}{\ensuremath{\theta}}
\newunicodechar{ρ}{\ensuremath{\rho}}
\newunicodechar{π}{\ensuremath{\pi}}
\newunicodechar{φ}{\ensuremath{\varphi}}
\newunicodechar{λ}{\ensuremath{\lambda}}
\newunicodechar{σ}{\ensuremath{\sigma}}
\newunicodechar{Δ}{\ensuremath{\Delta}}
\newunicodechar{Σ}{\ensuremath{\Sigma}}
\newunicodechar{ν}{\ensuremath{\nu}}
\newunicodechar{κ}{\ensuremath{\kappa}}
\newunicodechar{μ}{\ensuremath{\mu}}
\newunicodechar{τ}{\ensuremath{\tau}}
\newunicodechar{ω}{\ensuremath{\omega}}
\newunicodechar{ε}{\ensuremath{\varepsilon}}
\newunicodechar{Ψ}{\ensuremath{\Psi}}
\newunicodechar{⁻}{\ensuremath{{}^{-}}}
\newunicodechar{¹}{\ensuremath{{}^{1}}}
\newunicodechar{²}{\ensuremath{{}^{2}}}
\newunicodechar{³}{\ensuremath{{}^{3}}}
\newunicodechar{⋆}{\ensuremath{\star}}
\newunicodechar{∗}{\ensuremath{*}}
\newunicodechar{★}{\ensuremath{\bigstar}}
\newunicodechar{∎}{\ensuremath{\blacksquare}}
\newunicodechar{ℝ}{\ensuremath{\mathbb{R}}}
\newunicodechar{ℤ}{\ensuremath{\mathbb{Z}}}
\newunicodechar{₀}{\ensuremath{{}_{0}}}
\newunicodechar{₁}{\ensuremath{{}_{1}}}
\newunicodechar{₂}{\ensuremath{{}_{2}}}
\newunicodechar{₃}{\ensuremath{{}_{3}}}
\newunicodechar{₄}{\ensuremath{{}_{4}}}
\newunicodechar{₅}{\ensuremath{{}_{5}}}
\newunicodechar{₆}{\ensuremath{{}_{6}}}
\newunicodechar{₇}{\ensuremath{{}_{7}}}
\newunicodechar{₈}{\ensuremath{{}_{8}}}
\newunicodechar{₉}{\ensuremath{{}_{9}}}
\newunicodechar{✓}{\ensuremath{\checkmark}}
\newunicodechar{—}{---}
\newunicodechar{①}{(1)}
\newunicodechar{②}{(2)}
\newunicodechar{③}{(3)}
\newunicodechar{′}{\ensuremath{'}}
\newunicodechar{″}{\ensuremath{''}}
\newunicodechar{④}{(4)}
\newunicodechar{η}{\ensuremath{\eta}}
\newunicodechar{γ}{\ensuremath{\gamma}}
\newunicodechar{⇒}{\ensuremath{\Rightarrow}}
\newunicodechar{⟺}{\ensuremath{\Leftrightarrow}}
\newunicodechar{⇔}{\ensuremath{\Leftrightarrow}}
\usepackage{bookmark}
\IfFileExists{xurl.sty}{\usepackage{xurl}}{} % add URL line breaks if available
\urlstyle{same}
\hypersetup{
  hidelinks,
  pdfcreator={LaTeX via pandoc}}

\author{}
\date{}

\begin{document}

\section{The Generation Structure of Fermions: Particles as Completed
Closure Cycles in a Sea of
Light}\label{the-generation-structure-of-fermions-particles-as-completed-closure-cycles-in-a-sea-of-light}

\subsection{Spontaneous Differentiation of Forces from the Universal
Interaction, Spontaneous Differentiation of Particle Species by Parity,
Address, Coherence, and Hair, and the Counting of
Abundances}\label{spontaneous-differentiation-of-forces-from-the-universal-interaction-spontaneous-differentiation-of-particle-species-by-parity-address-coherence-and-hair-and-the-counting-of-abundances}

\textbf{Author:} Noriaki Kihara\\
\textbf{ORCID:} 0009-0004-6753-4020\\
\textbf{Date:} August 4, 2026\\
\textbf{Version DOI:} \texttt{10.5281/zenodo.21766707}\\
\textbf{Concept DOI:} \texttt{10.5281/zenodo.21766706}\\
\textbf{Position:} The Dimension Generation series, main paper 9. A
synthesis of the trilogy (two-grammar decomposition {[}3{]}, counting
readouts {[}4{]}, lock dynamics {[}5{]}) and the Paper-D preliminary
experiment suite {[}9{]}. \textbf{This paper is derivation and synthesis
only, and contains no new numerical experiments and no new figures.} All
quantitative facts rest on published papers and committed reproduction
packages, and every one of them can be traced one-to-one through the
source correspondence (Section 6).

\begin{center}\rule{0.5\linewidth}{0.5pt}\end{center}

\subsection{Claims --- The Mechanism by Which Particles Spontaneously
Differentiate from the Form of the
Wave}\label{claims-the-mechanism-by-which-particles-spontaneously-differentiate-from-the-form-of-the-wave}

The key is the inversion of default and exception.

\textbf{The default state is a sea of photons.} Zero-square-sum closure
(the fully coherent shared mode) is the axiom itself. Nothing needs to
be ``made'' --- the raw material of the system is lightlike from the
start.

\textbf{The first branching is parity.} Odd harmonics = fermionic; even
harmonics = bosonic.

\textbf{The second branching is address capture.} Within the coherent
sea, only those chains of exchange that happen to complete the closure
cycle of a register become recurring closures --- that is, particles.

\textbf{Matter and antimatter are given in pairs by the geometry of
cycle completion.} The halfway point of the cycle is \(\lambda^{62}=-1\)
(conjugate, opposite sign), and the full circuit is \(\lambda^{124}=+1\)
--- the sign of the double cover itself gives the pair.

\textbf{In other words, particles are rare events that completed the
closure cycle in the sea of light.} They are ``almost absent'' because
completing the closure cycle is difficult, and ``not zero'' because the
resolution is finite and the completion probability is strictly positive
--- both emerge automatically as structure.

What follows is nothing more than the derivation that led to this claim.

\begin{center}\rule{0.5\linewidth}{0.5pt}\end{center}

\subsection{Position of This Paper}\label{position-of-this-paper}

The paper-9 series has sought a mechanism by which both the species of
particles and the kinds of forces differentiate spontaneously from the
form of the wave alone, without either being given as input. The trilogy
established, at machine precision, the two-term identity of the
universal interaction and the unique differentiation of the two grammars
{[}3{]}, the structure of counting readouts {[}4{]}, and the
quantization mechanism together with the elimination approach to
selection {[}5{]}; the Paper-D preliminary experiment suite {[}9{]}
measured the register relation of the two alpha roots, clock blindness,
aliasing, and the mass beat. This paper synthesizes these into a single
generation structure. Every proof is either (i) a citation of a
published measurement or (ii) an elementary identity in the main text;
no new experiment is performed.

\subsection{1. The Problem}\label{the-problem}

Three points must be explained. (i) Where do the species of particles
come from --- without an external classification table. (ii) Where do
the kinds of forces come from --- without inputting coupling rules for
each particle species. (iii) Why is the abundance of matter ``extremely
small yet strictly nonzero'' --- without fine-tuning parameters.

\subsection{2. Derivation I --- The Universal Interaction: There Is Only
One
Force}\label{derivation-i-the-universal-interaction-there-is-only-one-force}

The interaction present in the system is one and only one, the
real-orthogonal exchange collision, and the transfer in a single
collision decomposes identically into the two terms

\[
dN_B=\underbrace{\sin^2\theta\,(N_A-N_B)}_{\text{magnitude term}}+\underbrace{\sin 2\theta\,\mathrm{Re}\langle a|b\rangle}_{\text{overlap term}}
\]

and \textbf{no third term exists} {[}3{]}. This identity does not ask
what the partner is. Hence this is the universal interaction, and the
distinction between forces is not a distinction between interactions but
\textbf{a distinction of which term the partner's waveform activates}:

\begin{itemize}
\tightlist
\item
  The magnitude term acts on everything that has a norm --- the gravity
  type (universal, additive, unscreenable, no quantization) {[}3{]}
\item
  The overlap term acts in the charge type only on pairs that share
  harmonics and carry hair --- the Coulomb type (amplitude product, ±,
  two neutralization-screening mechanisms, rational quantization)
  {[}3{]}
\item
  A fully shared relation carrying no hair bears neither charge and
  mediates both --- the photon (Section 3). Note that the two
  ``absences'' stem from separate axes: the absence of mass follows from
  full sharing (\(\det\Gamma=0\), Section 3, axis ③), and the absence of
  charge-type flow follows from the absence of hair (Section 3, axis ④)
\end{itemize}

This assignment is not a similarity of properties but a unique
correspondence that excludes the reverse assignment (Section 6.5 of
{[}3{]}). That is, \textbf{the kind of force is an alias of the kind of
particle, and the kind of particle is an alias of the form of the wave.}

\subsection{3. Derivation II --- The Four-Axis Decision Tree: Species
Differentiate from the Form of the
Wave}\label{derivation-ii-the-four-axis-decision-tree-species-differentiate-from-the-form-of-the-wave}

The species of a particle is determined solely by four independent
properties of the form of the wave.

\textbf{Axis ① Parity (statistics).} Odd harmonics are antiperiodic,
\(f(\theta+\pi)=-f(\theta)\), and two-valued --- fermionic, requiring a
double cover. Even harmonics are single-valued --- bosonic. The
fermionic construction from equal-amplitude bundles of odd harmonics and
their collisions have been measured {[}8{]}.

\textbf{Axis ② Binding (address and quantization).} A family of
harmonics bundled in integer ratios forms a harmonic family = a closure,
and recurs as an exact finite-order root of the exchange operator,
\(U^n=\pm I\) {[}2{]}. A closure carries a coprime address \(m/n\) on
the register \(n\), and \textbf{the address gives the charge value}: the
elementary charge is \(1-R=\sin^2(23\pi/124)=0.302822\), agreeing with
the dimensionless elementary charge \(\sqrt{4\pi\alpha}\) to 0.16 ppm
{[}2{]}{[}3{]}. The effective charge at high energy is a reading of the
same elementary charge on a fivefold-refined register
(\(620=5\times124\), clock \(1240=5\times248\)); to an observer with a
coarse clock, the fine root is read as exactly folded back onto the
electron's address (aliasing = a candidate mechanism for charge
universality) {[}9{]}. Relations that are not bound remain in the sea.

\textbf{Axis ③ Coherence (mass).} For a pair \((a,b)\), the Gram
invariant

\[
\det\Gamma=N_AN_B-|\langle a|b\rangle|^2
\]

satisfies \(\det\Gamma\ge0\) by Cauchy--Schwarz, with equality if and
only if \(b\propto a\) (full sharing) --- this is an identity, not an
experiment. Reading \(\det\Gamma\) as the mass-squared-type invariant,
\textbf{mass is incoherence, and for full sharing the reading called
mass can in principle never arise}.

\textbf{Axis ④ Hair (the carrier of charge).} Charge-type flow arises
only in shared pairs carrying ± carriers (hair), and neutralization
occurs through two mechanisms: zero components and channel non-sharing
{[}3{]}.

From the four axes the archetypes emerge:

\begin{itemize}
\tightlist
\item
  \textbf{The photon} = a fully shared even-difference mode carrying no
  hair. It is the quantum of the shared-channel selection rule
  \(\Delta k=\pm2\) (measured, coupling 0.5 {[}3{]}) --- the mode of the
  difference of two odd harmonics = even; a boson because it is even
  (axis ①), massless because it is fully shared (the identity of axis
  ③), chargeless because it carries no hair (axis ④), free because it
  holds no address. The four properties each follow from independent
  axes. The long-range inverse square is borne by the mediation of
  harmonic closure {[}7{]}
\item
  \textbf{The charged fermion (electron type)} = odd, bound (an address
  on register 124), incoherent bundle, shared with hair
\item
  \textbf{The neutral fermion} = odd, bound, hair-free or non-shared
  (subject to gravity only: the magnitude term reads the norm regardless
  of address {[}3{]})
\item
  \textbf{The baryon} = a bound closure topologically protected by
  winding number (baryon number = winding number; stability = the
  absence of a decay destination). Within this framework this
  identification is at the stage of a \textbf{correspondence hypothesis}
  (externally, the Skyrme lineage {[}12{]}{[}13{]} demonstrates an
  isomorphic structure), and its derivation is placed in Open Problem 4
\end{itemize}

One vacant slot remains explicit in the decision tree: \textbf{even,
bound, with hair} (a charged massive boson --- the slot for W±) --- Open
Problem 2.

\subsection{4. Derivation III --- Matter and Antimatter: The Geometry of
Cycle Completion Counts in
Pairs}\label{derivation-iii-matter-and-antimatter-the-geometry-of-cycle-completion-counts-in-pairs}

The eigenvalues of the elementary-charge register satisfy
\(\lambda^{31}=i,\ \lambda^{62}=-1,\ \lambda^{124}=+1\) {[}2{]}. The
halfway point of the completed cycle, \(\lambda^{62}=-1\), is
conjugation and sign reversal, and at the full circuit
\(\lambda^{124}=+1\) the system returns to itself. Combined with the 2:1
covering of the spinor (the state period is twice the observation period
{[}9{]}), \textbf{the geometry of completing the closure cycle itself
counts particle and antiparticle in pairs}. No additional symmetry is
needed to produce the pairs.

\subsection{5. Derivation IV --- Abundance: Scarcity and Nonzeroness
Emerge
Together}\label{derivation-iv-abundance-scarcity-and-nonzeroness-emerge-together}

A system of resolution \(N\) has \(M=N(N-1)/2\) relational waves
{[}6{]}. Only the integer-ratio pairs (pairs in which one is an integer
multiple of the other --- pairs bundled into a harmonic family, axis ②)
can complete the closure, and their number is

\[
\#\{(i,j):i<j,\ i\mid j\}=\sum_{i=1}^{N}\left(\left\lfloor\frac{N}{i}\right\rfloor-1\right)=N\ln N+(2\gamma-2)N+O(\sqrt N)
\]

(\(\gamma\): Euler's constant; Dirichlet's divisor-sum estimate). Hence
the fraction of relations that can become closures is

\[
f(N)=\frac{2\,(\ln N+2\gamma-2)}{N}+O(N^{-3/2})
\]

--- \textbf{it falls rapidly as \(N\) expands, yet never becomes zero}.
Both ``matter is almost absent'' and ``matter strictly exists'' emerge
automatically from this single fraction. Equating the observed
baryon-to-photon ratio \(\eta\approx6\times10^{-10}\) with \(f(N)\)
gives \(N\approx8\times10^{10}\) --- \textbf{η is read as a fossil of
the resolution at the instant the inflationary expansion halted in three
directions} (the value \(N\) itself is a reading from observation, not a
derivation --- Open Problem 8).

Baryogenesis in standard cosmology takes Sakharov's three conditions and
the magnitude of CP violation as external inputs {[}17{]}. The claim of
this framework is their replacement: \textbf{the smallness of the
asymmetry is not fine-tuning but the counting of integer-ratio pairs}
(the detailed correspondence with B-violating dynamics is left to the
Open Problems).

\subsection{6. Source Correspondence}\label{source-correspondence}

The verified facts on which each sentence of the Claims relies:

\begin{itemize}
\tightlist
\item
  \textbf{The sea of light = zero-square-sum closure}: the axiom system
  {[}1{]}. Full sharing ⇒ no mass is the Cauchy--Schwarz identity in the
  main text (Section 3, axis ③)
\item
  \textbf{The parity branching}: the measured fermionic construction and
  collisions of odd-harmonic bundles {[}8{]}, the implementation of ±
  hair {[}3{]}
\item
  \textbf{Address capture, recurrence, quantization}: exact finite-order
  roots {[}2{]}, counting readouts and Niven intersections {[}4{]},
  Arnold tongues, observational selection, and clock inheritance
  {[}5{]}, the register arithmetic of the two alpha roots, clock
  blindness, aliasing, and the mass beat {[}9{]}
\item
  \textbf{The λ periods of matter/antimatter}: Section 6.5 of {[}2{]}
\item
  \textbf{The accounting of abundance}: the mechanism of relational
  waves \(M=N(N-1)/2\) {[}6{]}; the counting is Section 5 of the main
  text (elementary)
\item
  \textbf{The universal interaction and force as an alias of species}:
  the two-term identity and the unique assignment {[}3{]}
\item
  \textbf{The photon = mediator}: the null theorem and the
  \(\Delta k=\pm2\) ladder {[}3{]}, the inverse-square mediation of
  harmonic closure {[}7{]}
\end{itemize}

\subsection{7. Relation to Prior Work --- Evidence of Convergence and
Differences}\label{relation-to-prior-work-evidence-of-convergence-and-differences}

The external literature is not the basis of the derivations in this
paper. It shows independent realizations of isomorphic configurations
and historical precedence.

\textbf{Particles as stable structures in a medium.} Kelvin's vortex
atoms {[}10{]}: atoms = stable vortex structures in a universal medium
--- the prototype of ``default = medium, particle = exceptional stable
structure.'' Wheeler's geons {[}11{]}: electromagnetic waves bundled by
their own gravity behaving like particles, ``mass without mass'' --- the
first quantitative attempt to build matter from light (difference:
requires a background spacetime and gravitational binding). Skyrme
{[}12{]} and Finkelstein--Rubinstein {[}13{]}: baryon number =
topological winding number, and spin statistics from the parity of the
winding number via the double cover --- the closest precedent for this
paper's baryon correspondence hypothesis and parity branching
(difference: solitons of fields on a background spacetime). Hestenes's
zitterbewegung electron {[}14{]}: the electron = a lightlike helical
closed orbit --- the single-particle version of ``mass = a bundle of
lightlike components'' (difference: no mechanism generating species).
Volovik {[}15{]}: fermionic quasiparticles, gauge fields, and effective
gravity co-emerge near the Fermi point of superfluid ³He --- a modern
realization of the ``co-emergence of species and forces'' (difference:
assumes a condensed-matter substructure, low-energy limit only, no
rational address quantization). Wilczek {[}16{]}: 99\% of the mass of
the visible universe comes from the dynamics of massless quarks and
gluons --- empirical support that ``mass from massless constituents'' is
the actual state of nature.

\textbf{Harmonics and particle species.} Kaluza--Klein {[}18{]}{[}19{]}:
the mode number on a compact circle gives charge (integer quantization)
and a mass tower --- the direct precedent for ``harmonic number on a
circle = charge'' (difference: requires extra spatial dimensions, and
the address is integer-only, with no coprime ratio, no parity structure,
and no universality mechanism). The RNS construction of string theory
{[}20{]}{[}21{]}: particle species = vibration modes of a single object,
and moreover periodic (Ramond) / antiperiodic (Neveu--Schwarz) boundary
conditions with integer / half-integer modes separate fermions from
bosons --- the direct precedent for the parity branching (difference:
requires a worldsheet, a critical dimension, and worldsheet
supersymmetry from outside; this paper uses only the parity of the
harmonics of the closed circle themselves). Regge--Chew--Frautschi
{[}22{]}: families of hadron species = ladders of excitations
(difference: phenomenology within the strong interaction). As a
historical seed, de Broglie {[}23{]}: matter = periodicity.

\textbf{Contrasting frame.} Sakharov {[}17{]}: the three conditions for
baryogenesis --- the standard frame whose replacement this paper's
counting claims.

Three things are unique to this paper: (i) statistics from the parity of
the harmonics themselves (no supersymmetry), (ii) charge from coprime
rational addresses (no extra dimensions), (iii) abundance from the
counting of finite resolution (no fine-tuning) --- and the fact that the
three branch simultaneously from a single axiom system.

\subsection{8. Open Problems}\label{open-problems}

\subsubsection{First tier --- gaps remaining directly in the claims of
this
paper}\label{first-tier-gaps-remaining-directly-in-the-claims-of-this-paper}

\begin{enumerate}
\def\labelenumi{\arabic{enumi}.}
\tightlist
\item
  \textbf{The address-selection problem (uniqueness of the elementary
  charge).} Register 124 has \(\varphi(124)=60\) coprime addresses, each
  giving a different charge value. At this stage the framework has no
  reason why only one of them is realized --- \textbf{there could just
  as well be 60 kinds of charge}. Eliminated: static counting {[}4{]},
  lock weights (width ratio \(>461\)) {[}5{]}, fixed points {[}5{]},
  observational selection (inherited into the clock) {[}5{]}, four
  simple arithmetic screenings {[}9{]}. Remaining candidates: mass
  balance (norm-sector dynamics, untested) and clock inheritance from
  the parent closure. The final form is ``why is the observation clock
  commensurable with 248?'' {[}5{]}.
\item
  \textbf{The boson spectrum problem.} The only boson derived is the
  photon. The reason for the existence of massive bosons and charged
  bosons (the W± slot = even, bound, with hair) has not been derived; it
  is the vacant slot in the decision tree.
\item
  \textbf{The mass-dictionary problem (energy hierarchy).} Mass =
  incoherence is qualitatively established, but the functional form
  giving physical mass from register and address is underived. The naive
  candidate (register ratio 5, coupling ratio 25) is inconsistent with
  the generation ratio \(m_\mu/m_e=206.77\) by the model's own precision
  standard (rejected {[}9{]}). Why the energy hierarchy appears as
  discrete generations is also underived. \textbf{We can state the
  existence of massive things, but we cannot state a single mass value.}
\item
  \textbf{The baryon identification problem.} Baryon number = winding
  number is at the stage of a correspondence hypothesis; its derivation
  within this framework (the definition of the winding number, its
  conservation, and stability = the absence of a destination) is
  incomplete.
\end{enumerate}

\subsubsection{Second tier --- connection problems shared with adjacent
papers}\label{second-tier-connection-problems-shared-with-adjacent-papers}

\begin{enumerate}
\def\labelenumi{\arabic{enumi}.}
\setcounter{enumi}{4}
\tightlist
\item
  \textbf{The distance dictionary}: the translation
  \(\Delta\theta\leftrightarrow r\) and the Kepler-type closure
  condition \(\omega^2R''^3=\text{const.}\) (shared with the trilogy
  {[}3{]} and the real-center paper {[}7{]}).
\item
  \textbf{Sign translation}: whether the ± of the relative phase maps to
  spatial attraction or repulsion (already fixed as a prediction that
  will be decided automatically the moment the dictionary is completed
  {[}3{]}).
\item
  \textbf{Geometric connection}: the eight problems of the realization
  map from three directions to 3+1-dimensional spacetime are held by the
  closing section of Part C of the trilogy {[}5{]}. This paper does not
  duplicate them.
\end{enumerate}

\subsubsection{Third tier --- underived items remaining on the side of
the
premises}\label{third-tier-underived-items-remaining-on-the-side-of-the-premises}

\begin{enumerate}
\def\labelenumi{\arabic{enumi}.}
\setcounter{enumi}{7}
\tightlist
\item
  \textbf{The value of the resolution \(N\)}: \(f(N)\) is a derivation,
  but \(N\approx10^{11}\) at the halt of expansion is a reading from the
  observed ratio \(\eta\), not a derivation. Why the expansion halts at
  that \(N\) and in three directions is identical to the open problem of
  {[}6{]}.
\item
  \textbf{The place of the strong and weak interactions}: the menu of
  forces in this framework consists only of the two grammars and the
  mediator; SU(2)- and SU(3)-type structures are untouched.
\item
  \textbf{The value of spin}: the parity ↔ statistics correspondence is
  established, but from which discrete feature of the internal modes the
  value of spin as a transformation factor (\(\pm1/2,\ \pm1\)) is read
  out is underived.
\end{enumerate}

This table of problems is not an enumeration of the weaknesses of this
paper but a guide for further progress of the series.

\subsection{9. Reproducibility}\label{reproducibility}

This paper is derivation and synthesis only and contains no new
numerical experiments. All the quantitative facts it relies on are
already published: the reproduction packages of the trilogy
{[}3{]}{[}4{]}{[}5{]} (every number in one-to-one correspondence with
committed runners, with SHA-256 verification), the reproduction bundle
of the finite-order roots {[}2{]}, the reproduction package of the
N-body relational waves {[}6{]}, the derivations and verifications of
harmonic closure and the real center {[}7{]}, and the Paper-D
preliminary experiment suite {[}9{]} (in the repository, commits
7beea7ed and d802907a, with predictions written in advance and a full
record of refutations and corrections). The only new mathematics used in
the main text consists of two elementary facts: the equality condition
of Cauchy--Schwarz (Section 3) and Dirichlet's divisor-sum estimate
(Section 5).

\begin{center}\rule{0.5\linewidth}{0.5pt}\end{center}

\section{References}\label{references}

\subsection{Self-citations}\label{self-citations}

\begin{enumerate}
\def\labelenumi{\arabic{enumi}.}
\tightlist
\item
  Noriaki Kihara, ``Anonymous Equal-Amplitude Composite-Wave Model:
  Basic Axiom System'', Concept DOI: \texttt{10.5281/zenodo.21315735}
  (always the latest version), 2026.
\item
  Noriaki Kihara, ``Discovery of Finite-Order Resonance in Iterated
  Exchange Scattering --- Identifying the Origin of Peaks near
  Fine-Structure-Constant Values 137 and 128 with a Reproducible
  Wave-Packet Model'', Version DOI: \texttt{10.5281/zenodo.21421367},
  Concept DOI: \texttt{10.5281/zenodo.21421366}, 2026.
\item
  Noriaki Kihara, ``Two-Grammar Decomposition of Interaction in an
  Anonymous Two-Channel Closed Wave System'' (Part I of the trilogy),
  Version DOI: \texttt{10.5281/zenodo.21763996}, Concept DOI:
  \texttt{10.5281/zenodo.21763995}, 2026.
\item
  Noriaki Kihara, ``Structure of Counting Readouts in an Anonymous
  Two-Channel Closed Wave System'' (Part II of the trilogy), Version
  DOI: \texttt{10.5281/zenodo.21763998}, Concept DOI:
  \texttt{10.5281/zenodo.21763997}, 2026.
\item
  Noriaki Kihara, ``Lock Dynamics in an Anonymous Two-Channel Closed
  Wave System'' (Part III of the trilogy), Version DOI:
  \texttt{10.5281/zenodo.21764000}, Concept DOI:
  \texttt{10.5281/zenodo.21763999}, 2026.
\item
  Noriaki Kihara, ``Arrest of Spontaneous Splitting and the Creation of
  an Additional Axis in an N-Body Relational-Wave Closed System'',
  Version DOI: \texttt{10.5281/zenodo.21543071}, Concept DOI:
  \texttt{10.5281/zenodo.21543070}, 2026; and ``Temporal Structure of
  Three-Direction Generation in an N-Body Relational-Wave Closed System
  --- Causal Separation by Two-Stage Seed Removal'', Version DOI:
  \texttt{10.5281/zenodo.21614403}, Concept DOI:
  \texttt{10.5281/zenodo.21614402}, 2026. (Relational waves
  \(M=N(N-1)/2\), inflationary expansion, three-direction halt, seed
  independence)
\item
  Noriaki Kihara, ``Future Phase-Position Acceleration Map and the
  Inverse-Square Law via Harmonic Closure in an AB Two-Body Closed Phase
  System v4'', Concept DOI: \texttt{10.5281/zenodo.21441081}; and
  ``Rederivation of the Acceleration Map by a Real-Center Reading in a
  Closed Two-Body AB Phase System'', Version DOI:
  \texttt{10.5281/zenodo.21765368}, Concept DOI:
  \texttt{10.5281/zenodo.21765367}, 2026.
\item
  Noriaki Kihara, ``Preliminary Summary of Acceleration-Basis and
  Localization Exchange in Exchange-Interference Scattering-Matrix
  Fermion-Like Collisions v1'', Version DOI:
  \texttt{10.5281/zenodo.21333768}, Concept DOI:
  \texttt{10.5281/zenodo.21333766}, 2026. (The measured fermionic
  construction and collisions of the odd-harmonic bundle B63)
\item
  Noriaki Kihara, Paper-D preliminary experiment suite (register
  arithmetic of the two alpha roots, clock blindness, aliasing / mass
  beat / address-selection survey), \texttt{paperD\_*} experiment
  folders in the repository \texttt{ai-chat-logs-open}, commits
  \texttt{7beea7ed} and \texttt{d802907a}, 2026.
\end{enumerate}

\subsection{External References}\label{external-references}

\begin{enumerate}
\def\labelenumi{\arabic{enumi}.}
\setcounter{enumi}{9}
\tightlist
\item
  W. Thomson (Lord Kelvin), ``On Vortex Atoms'', \emph{Proc. R. Soc.
  Edinburgh} \textbf{6}, 94--105 (1867).
\item
  J. A. Wheeler, ``Geons'', \emph{Phys. Rev.} \textbf{97}, 511 (1955).
  DOI: \texttt{10.1103/PhysRev.97.511}.
\item
  T. H. R. Skyrme, ``A non-linear field theory'', \emph{Proc. R. Soc.
  Lond. A} \textbf{260}, 127 (1961); ``A unified field theory of mesons
  and baryons'', \emph{Nucl. Phys.} \textbf{31}, 556 (1962).
\item
  D. Finkelstein and J. Rubinstein, ``Connection between Spin,
  Statistics, and Kinks'', \emph{J. Math. Phys.} \textbf{9}, 1762
  (1968). DOI: \texttt{10.1063/1.1664510}.
\item
  D. Hestenes, ``The zitterbewegung interpretation of quantum
  mechanics'', \emph{Found. Phys.} \textbf{20}, 1213 (1990). DOI:
  \texttt{10.1007/BF01889466}.
\item
  G. E. Volovik, \emph{The Universe in a Helium Droplet}, Oxford
  University Press (2003).
\item
  F. Wilczek, ``Mass without mass I: Most of matter'', \emph{Phys.
  Today} \textbf{52}(11), 11 (1999); ``Mass without mass II: The medium
  is the mass-age'', \emph{Phys. Today} \textbf{53}(1), 13 (2000).
\item
  A. D. Sakharov, ``Violation of CP invariance, C asymmetry, and baryon
  asymmetry of the universe'', \emph{JETP Lett.} \textbf{5}, 24 (1967).
\item
  T. Kaluza, ``Zum Unitätsproblem der Physik'', \emph{Sitzungsber.
  Preuss. Akad. Wiss. Berlin}, 966 (1921).
\item
  O. Klein, ``Quantentheorie und fünfdimensionale Relativitätstheorie'',
  \emph{Z. Phys.} \textbf{37}, 895 (1926). DOI:
  \texttt{10.1007/BF01397481}.
\item
  P. Ramond, ``Dual Theory for Free Fermions'', \emph{Phys. Rev.~D}
  \textbf{3}, 2415 (1971). DOI: \texttt{10.1103/PhysRevD.3.2415}.
\item
  A. Neveu and J. H. Schwarz, ``Factorizable dual model of pions'',
  \emph{Nucl. Phys. B} \textbf{31}, 86 (1971).
\item
  G. F. Chew and S. C. Frautschi, ``Principle of Equivalence for All
  Strongly Interacting Particles within the S-Matrix Framework'',
  \emph{Phys. Rev.~Lett.} \textbf{7}, 394 (1961). DOI:
  \texttt{10.1103/PhysRevLett.7.394}.
\item
  L. de Broglie, ``Recherches sur la théorie des quanta'', \emph{Ann.
  Phys.} (Paris) \textbf{3}, 22 (1925) (1924 doctoral thesis).
\end{enumerate}

\begin{center}\rule{0.5\linewidth}{0.5pt}\end{center}

\textbf{Addendum (the evidence structure of this paper).} Every
quantitative claim of this paper corresponds one-to-one to one of the
following: (i) measurements in published papers (verifiable through the
reproduction package of each DOI), (ii) committed preliminary
experiments in the repository (with predictions written in advance and a
full record of refutations and corrections), or (iii) elementary
identities in the main text. It contains no new experiments and no new
figures because the claim of this paper lies not in new data but in the
\textbf{arrangement} of established facts --- the inversion of default
and exception.

\end{document}
